problem:  
Let \( M^5 \) be a closed, simply connected Riemannian manifold with **nonnegative sectional curvature** and an **effective, isometric \(T^2\)-action**, i.e. symmetry rank \(2\) (one less than maximal).  Write \( X = M/T^2 \) for the orbit space, which is a compact, 3‑dimensional Alexandrov space with curvature \(\ge0\).  

Assume the action is **almost maximal** in the sense that the \(T^2\)-fixed point set is nonempty and the principal orbits are dense.  Such an action has a one‑dimensional singular set in \(X\) consisting of images of singular orbits with circle isotropy, and isolated points corresponding to fixed orbits.  

**Open problem:**  
Does every such \( (M^5, g) \) appear as an **equivariant Gromov–Hausdorff limit** of 5‑manifolds \( (M_i, g_i) \) with **nonnegative curvature** and **effective isometric \(T^3\)-actions** (of maximal symmetry rank), where one circle factor in the \(T^3\) action collapses uniformly?  

Equivalently, is every nonnegatively curved \(5\)-manifold with an effective \(T^2\)-isometric action a **curvature‑preserving collapse** of a manifold with one more commuting symmetry?  
If not, characterize which orbit‑space topologies \(X=M/T^2\) are *non‑degenerable*, i.e. cannot arise as a limit of convex polyhedral \(2\)-ball orbit spaces coming from \(T^3\)-actions.  

Why is it a "good" problem:  
- **Profound insight:** The question probes whether curvature and symmetry interact *continuously* at the jump from maximal to almost maximal rank.  If smooth \(T^2\)-actions can always be realized as curvature‑preserving degenerations of \(T^3\)-actions, it would reveal an unforeseen **stability phenomenon** for symmetry under nonnegative curvature.  
- **Bridge between fields:** The problem lies precisely at the crossroads of **Riemannian transformation groups**, **Alexandrov geometry**, and **metric collapse theory** (in the sense of Cheeger–Gromov–Fukaya).  Progress would necessarily merge ideas from geometric analysis (equivariant Gromov–Hausdorff limits) with the topological classification of torus actions.  
- **Extreme simplicity:**  One can phrase it succinctly as:  
  “Is every nonnegatively curved \(5\)-manifold with an isometric \(T^2\)-action a curvature‑preserving collapse limit of a \(T^3\)-symmetric one?”  
  The statement is simple, yet resolving it would likely require novel insights into how curvature constraints control the evolution of symmetry, a deep rigidity problem bridging metric geometry and topology.